\section{The Mal'cev condition and the universal comparison}
\label{d:begin}

Suppose a relation connects inputs in $X$ with outputs in $Y$. Three
corners of a rectangle may already force the fourth:
\[
x\,R\,y,\qquad x'\,R\,y,\qquad x'\,R\,y'
\quad\Longrightarrow\quad x\,R\,y'.
\]
This property is called \emph{difunctionality}. It says that the
relation breaks into complete rectangular blocks. It does not say that
each input has just one output. Several inputs and several outputs can
belong to the same block.

The structural comparison is already suggested by this picture. Group
the inputs according to which block they occupy, and group the outputs
in the same way. The two resulting quotient objects should agree. We
will construct that agreement, including its maps, before using it to
compare any filtrations.

\subsection{Why relative difference is enough}

\begin{definition}[Mal'cev category]\label{d:malcev}
A regular category is a \emph{Mal'cev category} if every reflexive
internal relation is an equivalence relation.
\end{definition}

This categorical condition has an elementary source. Suppose we have
algebraic structures with a ternary term $p$ satisfying
\[
p(x,y,y)=x,\qquad p(x,x,y)=y.
\]
If a compatible relation $R\subseteq X\times X$ is reflexive and
$x\,R\,y$, its closure under $p$ gives
\[
p\bigl((y,y),(x,y),(x,x)\bigr)=(y,x)\in R.
\]
Thus it is symmetric. If also $y\,R\,z$, closure gives
\[
p\bigl((x,y),(y,y),(y,z)\bigr)=(x,z)\in R,
\]
so it is transitive. The operation is applied coordinatewise to pairs.
For groups $p(x,y,z)=xy^{-1}z$; for modules it is $x-y+z$. No
commutativity of group multiplication was used. A term is an operation
on the underlying algebra; in the group case it need not be a group
homomorphism from $X^3$ to $X$.

This gives the central common resource: a way to take a relative
difference and then reattach it at another point. The categorical
definition retains its effect on compatible relations without requiring
a chosen formula for the operation.

\begin{proposition}[Three forms of the Mal'cev condition]\label{d:difunctional}
For a regular category the following conditions are equivalent:
\begin{enumerate}
\item every reflexive relation is an equivalence relation;
\item equivalence relations on the same object commute under
relation composition;
\item every relation $R:X\rightsquigarrow Y$ is difunctional, that is,
\[
RR^{\mathrm{op}}R=R.
\]
\end{enumerate}
\end{proposition}
\begin{proof}
Assume the first condition. If $E,F$ are equivalence relations on $X$,
their composite $EF$ is reflexive, since both contain $\Delta_X$.
It is therefore symmetric, and
\[
EF=(EF)^{\mathrm{op}}=FE.
\]
This proves the second condition.

Assume the second. Give $R$ a carrier with projections
$a:M\to X$ and $b:M\to Y$. The kernel pairs
\[
E_a=a^{\mathrm{op}}a,\qquad E_b=b^{\mathrm{op}}b
\]
are equivalence relations on the \emph{same} object $M$, so they
commute. Using Lemma~\ref{f:graph-kernel}, we obtain
\begin{align*}
RR^{\mathrm{op}}R
&=ba^{\mathrm{op}}ab^{\mathrm{op}}ba^{\mathrm{op}}\\
&=bE_aE_ba^{\mathrm{op}}
 =bE_bE_aa^{\mathrm{op}}\\
&=(bb^{\mathrm{op}}b)(a^{\mathrm{op}}aa^{\mathrm{op}})
 =ba^{\mathrm{op}}=R.
\end{align*}
The equalities are typed on $M$; we have not interchanged unrelated
relations on $X$ and $Y$. This proves the third condition.

Finally assume the third and let $R$ be reflexive. Monotonicity gives
\[
R^{\mathrm{op}}=\Delta_XR^{\mathrm{op}}\Delta_X
\leq RR^{\mathrm{op}}R=R.
\]
Taking opposites gives the reverse inclusion, so $R$ is symmetric.
Also $\Delta_X\leq R^{\mathrm{op}}$, and hence
$RR\leq RR^{\mathrm{op}}R=R$. Thus $R$ is transitive and is an
equivalence relation.
\end{proof}

The local-witness interpretation of $RR^{\mathrm{op}}R\leq R$ is
exactly the three-corners implication above. The reverse inclusion
always holds: a related pair can supply all three edges of a degenerate
rectangle. The proposition explains why the rectangle-completion law
holds simultaneously for every relation, rather than having to be
verified anew for each quotient comparison.

\subsection{Partial identifications and their supports}

Relations used in a filtration may live on a subobject rather than all
of the ambient object. We must record that support instead of silently
assuming reflexivity everywhere.

\begin{lemma}[Partial equivalence relations]\label{d:partial-equivalence}
Let $E:X\rightsquigarrow X$ be symmetric and transitive in a regular
category. The two projections of its carrier have the same image
$m:D\hookrightarrow X$. Restricted to $D$, the relation is an
equivalence relation $E_D$. In a Barr-exact category there is a regular
epimorphism $q:D\twoheadrightarrow Q$ with
\[
E=mq^{\mathrm{op}}qm^{\mathrm{op}}.
\]
Conversely, every relation of this displayed form is symmetric and
transitive, with support $D$.
\end{lemma}
\begin{proof}
Symmetry interchanges the two projections, so their images agree.
Their factorizations through $D$ give a relation $E_D\hookrightarrow
D\times D$; the pair is still jointly monic. Both of its projections
are regular epic. If $x$ is an input in $D$, a cover supplies a witness
$x\,E_D\,y$. Symmetry gives $y\,E_D\,x$, and transitivity gives
$x\,E_D\,x$. Lemma~\ref{f:local-calculus} descends this membership,
so $E_D$ is reflexive. Symmetry and transitivity are inherited by
restriction. Exactness gives $E_D=q^{\mathrm{op}}q$, proving the formula.

Conversely, $q^{\mathrm{op}}q$ is an equivalence relation on $D$.
Its extension along $m$ is symmetric; since $m^{\mathrm{op}}m=\Delta_D$,
its square is itself. The projections onto $D$ are regular epic
because the relation is reflexive, so their images in $X$ are $D$.
\end{proof}

We call $Q$ the quotient of the partial equivalence relation. It is
the quotient of its support, not a quotient of all of $X$. This
distinction permits the same relation language to describe every step
of a chain whose terms are only normal in their successors.

For example, a quotient map $q:A\twoheadrightarrow S$ with
$A\hookrightarrow G$ determines a relation on $G$: two inputs are
related precisely when they both lie in $A$ and have the same image
in $S$. Its support is $A$; in a pointed category its zero class is
$\Ker(q)$. Thus one partial relation remembers the numerator, the
identifications, and the chosen quotient, up to its unique quotient
isomorphism.

\subsection{One relation, two quotients, one object}

Let $R:X\rightsquigarrow Y$ have carrier $M$ and factor its projections
through their images:
\[
\begin{tikzcd}[column sep=large]
&M \arrow[dl,two heads,"a"'] \arrow[dr,two heads,"b"]&\\
X' \arrow[d,hook,"m"'] &&Y' \arrow[d,hook,"n"]\\
X&&Y.
\end{tikzcd}
\]
The pair $(a,b)$ is jointly monic. We use $\widehat R:X'\rightsquigarrow Y'$
for the relation on these supports, so $R=n\widehat Rm^{\mathrm{op}}$.
Both projections of $\widehat R$ are regular epic.

\begin{theorem}[The universal support-quotient comparison]\label{d:comparison}
In a Barr-exact Mal'cev category, put
\[
E=\widehat R^{\mathrm{op}}\widehat R
\quad\text{on }X',\qquad
F=\widehat R\widehat R^{\mathrm{op}}
\quad\text{on }Y'.
\]
These are equivalence relations. For their regular quotient maps
$q:X'\twoheadrightarrow Q=X'/E$ and
$p:Y'\twoheadrightarrow P=Y'/F$, there is a unique isomorphism
$\theta:Q\iso P$ satisfying
\[
\theta qa=pb.
\]
In particular the two support quotients of any relation are canonically
isomorphic.
\end{theorem}
\begin{proof}
Because $a$ is regular epic, every input in $X'$ has locally some
$\widehat R$-related output. Using that same output twice proves
$\Delta_{X'}\leq E$. Likewise $\Delta_{Y'}\leq F$. Both relations are
symmetric by taking opposites. Difunctionality gives
\[
E^2=\widehat R^{\mathrm{op}}
 (\widehat R\widehat R^{\mathrm{op}}\widehat R)=E,
\qquad
F^2=(\widehat R\widehat R^{\mathrm{op}}\widehat R)
 \widehat R^{\mathrm{op}}=F.
\]
Thus they are equivalence relations, and exactness supplies $q,p$.

We construct the map rather than merely matching equivalence classes.
The map $pb:M\to P$ is equal on the kernel pair of $a$: if $au=av$,
then $bu$ and $bv$ share a $\widehat R$-related input, so they are
$F$-related and $pbu=pbv$. Lemma~\ref{f:descent} gives a unique
$g:X'\to P$ with $ga=pb$.

Next $g$ is equal on $E$. An $E$-related pair is locally of the form
$a(u),a(v)$ with $b(u)=b(v)$. Consequently
$ga(u)=pb(u)=pb(v)=ga(v)$. These equalities hold after a cover,
and covers are epimorphisms, so they hold on the relation itself.
Hence $g$ descends uniquely through $q$, giving $\theta:Q\to P$.

The same construction with the opposite relation gives $\psi:P\to Q$
such that $\psi pb=qa$. It follows that
\[
\psi\theta qa=qa,\qquad \theta\psi pb=pb.
\]
Both $qa$ and $pb$ are composites of regular epimorphisms, hence are
epic. Cancellation proves $\psi\theta=\id_Q$ and
$\theta\psi=\id_P$. Finally $qa$ is epic, so the equation specifying
$\theta$ proves its uniqueness.
\end{proof}

The construction is summarized by the commutative diagram
\[
\begin{tikzcd}[column sep=large,row sep=large]
&M \arrow[dl,two heads,"a"'] \arrow[dr,two heads,"b"]&\\
X' \arrow[d,two heads,"q"'] &&Y' \arrow[d,two heads,"p"]\\
Q \arrow[rr,"\theta", "\sim"'] &&P.
\end{tikzcd}
\]
There is generally no map $X'\to Y'$ to put across the middle row.
The relation supplies multiple possible outputs. The quotient removes
precisely their ambiguity, after which the comparison becomes a
uniquely determined isomorphism.

For a familiar model, take the additive group $\mathbb Z$ and the relation
$xRy$ meaning that $x-y$ is even. This is a subgroup of
$\mathbb Z\times\mathbb Z$, with both supports equal to $\mathbb Z$.
Two inputs share an output precisely when they have the same parity;
the same is true of two outputs sharing an input. Thus both support
quotients are $\mathbb Z/2\mathbb Z$, and the comparison sends the parity
class of an input to that of any related output. There are infinitely
many possible outputs before quotienting, but exactly one class afterwards.
This illustrates the construction; the categorical theorem did not need
a preferred output or a choice of representatives.

Here is another view of the same map, useful for computations. Its
graph is
\[
L=p\widehat Rq^{\mathrm{op}}:Q\rightsquigarrow P.
\]
Indeed,
\begin{align*}
L^{\mathrm{op}}L
 &=q\widehat R^{\mathrm{op}}p^{\mathrm{op}}p
       \widehat Rq^{\mathrm{op}}
   =qE^2q^{\mathrm{op}}=\Delta_Q,\\
LL^{\mathrm{op}}
 &=p\widehat Rq^{\mathrm{op}}q
       \widehat R^{\mathrm{op}}p^{\mathrm{op}}
   =pF^2p^{\mathrm{op}}=\Delta_P.
\end{align*}
We used $q^{\mathrm{op}}q=E$, $p^{\mathrm{op}}p=F$, their
idempotence, and $qq^{\mathrm{op}}=\Delta_Q$,
$pp^{\mathrm{op}}=\Delta_P$. Lemma~\ref{f:functional} makes $L$ a
map, and makes $L^{\mathrm{op}}$ its inverse. This is the same $\theta$
because its graph relates the quotient images of every pair in
$\widehat R$.

This theorem is the comparison mechanism from which the refinement
rectangle will be built. It is stronger than a statement about lengths:
it constructs the actual isomorphism between the surviving objects.
No filtration, simplicity assumption, or Jordan--H\"older theorem has
entered its proof.

\subsection{Which relations descend through a quotient?}

There is one further compatibility to make explicit. If we identify
inputs by a quotient map, a relation descends only when changing an
input within its equivalence class does not change the relation. This
is saturation.

\begin{proposition}[Saturation and descent]\label{d:saturated-descent}
Let $q:X\twoheadrightarrow Q$ and $p:Y\twoheadrightarrow P$ be
regular epimorphisms, with kernel pairs
$E=q^{\mathrm{op}}q$ and $F=p^{\mathrm{op}}p$. There is an
order-preserving bijection
\[
\left\{\text{relations }L:Q\rightsquigarrow P\right\}
\longleftrightarrow
\left\{R:X\rightsquigarrow Y\mid FRE=R\right\}.
\]
It sends $L$ to $p^{\mathrm{op}}Lq$, and its inverse sends $R$ to
$pRq^{\mathrm{op}}$. The first construction is precisely inverse image
of the relation along $q\times p$.
\end{proposition}
\begin{proof}
Inverse image consists of pairs $(x,y)$ whose quotient images satisfy
$q(x)\,L\,p(y)$, which is $p^{\mathrm{op}}Lq$ by the graph formulas.
Since $qq^{\mathrm{op}}=\Delta_Q$ and
$pp^{\mathrm{op}}=\Delta_P$, the two constructions satisfy
\[
p(p^{\mathrm{op}}Lq)q^{\mathrm{op}}=L,
\qquad
p^{\mathrm{op}}(pRq^{\mathrm{op}})q=FRE.
\]
An inverse-image relation is saturated because
$Fp^{\mathrm{op}}=p^{\mathrm{op}}$ and $qE=q$. Thus the formulas are
mutually inverse exactly on the stated saturated relations. Both are
order preserving by monotonicity of relation composition.
\end{proof}

For a relation $T$ on $X$ and a single quotient $q$, saturation is
$ETE=T$. Since $E$ contains the diagonal, this implies separately
$ET=T$ and $TE=T$. Consequently the correspondence respects
composition of saturated relations:
\[
(qSq^{\mathrm{op}})(qTq^{\mathrm{op}})
=qSETq^{\mathrm{op}}=qSTq^{\mathrm{op}}.
\]
It also respects opposites. Thus a saturated symmetric transitive
relation descends to a symmetric transitive relation, including its
possibly proper support. Conversely, inverse image preserves these
properties: for example
\[
(q^{\mathrm{op}}Lq)^2
=q^{\mathrm{op}}L^2q\leq q^{\mathrm{op}}Lq
\]
when $L$ is transitive. These statements hold for partial equivalence
relations as well as for equivalence relations on all of $X$.

\begin{corollary}[Quotients of quotients]\label{d:iterated-quotient}
In a Barr-exact category, fix $q:X\twoheadrightarrow Q$ with kernel
pair $E$. Equivalence relations $L$ on $Q$ correspond to equivalence
relations $T$ on $X$ containing $E$, by
\[
T=q^{\mathrm{op}}Lq,\qquad L=qTq^{\mathrm{op}}.
\]
If $t:X\twoheadrightarrow Z$ is the quotient of $T$, it factors
uniquely as $t=uq$, where $u:Q\twoheadrightarrow Z$ is the quotient
of $L$. Thus $(X/E)/L\cong X/T$, with the isomorphism determined by
the quotient maps.
\end{corollary}
\begin{proof}
If $T$ is an equivalence relation containing $E$, then
$T\leq ETE\leq TTT=T$, so it is saturated. Its descent is reflexive
because $q\Delta_Xq^{\mathrm{op}}=\Delta_Q$, and is symmetric and
transitive by the preceding formulas. Conversely, the inverse image
of an equivalence relation is an equivalence relation containing
$q^{\mathrm{op}}\Delta_Qq=E$.

Since $E\leq T$, the map $t$ is equal on the kernel pair of $q$,
so it descends uniquely to $u$. Factor $u$ into its image. If its mono
part were a proper subobject of $Z$, then $t=uq$ would factor through
that same subobject, contradicting the lifting property of the regular
epimorphism $t$. Hence $u$ is regular epic. Its kernel pair is
\[
u^{\mathrm{op}}u
=qq^{\mathrm{op}}u^{\mathrm{op}}uqq^{\mathrm{op}}
=qt^{\mathrm{op}}tq^{\mathrm{op}}=qTq^{\mathrm{op}}=L.
\]
Lemma~\ref{f:descent} makes it the quotient of $L$. The universal
property proves the claimed canonical isomorphism.
\end{proof}

This last result needs exactness and regularity, but not the Mal'cev
condition. Recording that distinction matters: Mal'cev is used to make
\emph{every} relation difunctional and hence to produce the universal
comparison. The ordinary law for compatible iterated quotients has
already been supplied by exactness. Later, in a higher setting, we
will see why retaining such compatibility data remains essential.

The relation calculus and support comparison are classical. For
context, see \href{https://arxiv.org/pdf/math/0610552}{Knop,
\emph{Tensor envelopes of regular categories}}, Section~5 and
Appendices~A--B. The arguments here have been written out so that the
refinement theorem can depend on explicit constructions rather than on
an unnamed isomorphism theorem.
